\subsection{On known pockets the Jev-like networks keep most of the recall of the joint network}

The networks were trained with the same recipe on the docking scores of known pockets, and every network was evaluated on the pockets of its training data against 10{,}000 ligands per pocket that were not used in training (Table~\ref{tab:known}). The checkpoint of the first group of networks was chosen on seen pockets, as the setting of this study asks, in five cross-fitted folds that give every pocket four networks per condition, and the evaluation ligands were not used in the choice of these checkpoints. Within a budget of the best 10\% of the ranking, the dual encoder recovers \kfDualTen{} of the top-1\% docking hits of the \kfNTargets{} pockets (95\% interval \kfDualTenLo{} to \kfDualTenHi{}), the pooled network \kfPooledTen{} (\kfPooledTenLo{} to \kfPooledTenHi{}) and the cross-attention network \kfXattnTen{} (\kfXattnTenLo{} to \kfXattnTenHi{}). A network that has one constant dummy pocket for every target recovers \kfConstTen{} (\kfConstTenLo{} to \kfConstTenHi{}), the recall that the ligand supports alone, so that the pocket state adds \kfDiffDualConstTen{} (\kfDiffDualConstTenLo{} to \kfDiffDualConstTenHi{}) for the dual encoder, \kfDiffPooledPooledConstTen{} (\kfDiffPooledPooledConstTenLo{} to \kfDiffPooledPooledConstTenHi{}) for the pooled network and \kfDiffXattnXattnConstTen{} (\kfDiffXattnXattnConstTenLo{} to \kfDiffXattnXattnConstTenHi{}) for the cross-attention network.

The cross-attention network is the most accurate of the three. Its recall exceeds that of the dual encoder by \kfDiffXattnDualTen{} (\kfDiffXattnDualTenLo{} to \kfDiffXattnDualTenHi{}) and that of the pooled network by \kfDiffXattnPooledTen{} (\kfDiffXattnPooledTenLo{} to \kfDiffXattnPooledTenHi{}), so that the dual encoder keeps \kfRatioDual{}\% (\kfRatioDualLo{} to \kfRatioDualHi{}) and the pooled network \kfRatioPooled{}\% (\kfRatioPooledLo{} to \kfRatioPooledHi{}) of its recall at the 10\% budget. The pooled network has the cost of the dual encoder and recovers \kfDiffPooledDualTen{} (\kfDiffPooledDualTenLo{} to \kfDiffPooledDualTenHi{}) more of the hits. At the 20\% budget the three networks recover \kfDualTwenty{}, \kfPooledTwenty{} and \kfXattnTwenty{}. A longer recipe (learning rate $10^{-4}$, up to 150{,}000 updates) raises the recall of the dual encoder to \kfDualLongTen{} and that of the cross-attention network to \kfXattnLongTen{}, and the dual encoder keeps \kfRatioDualLong{}\% (\kfRatioDualLongLo{} to \kfRatioDualLongHi{}) of the recall of the cross-attention network. The price of the additional accuracy is the cost of the previous section: the model time of the cross-attention network for 57 pockets is \tmARatioXaHalfFiftySevenModel{} to \tmARatioXaSingleFiftySevenModel{} times that of the dual encoder, and \tmARatioXaHalfThousandModel{} to \tmARatioXaSingleThousandModel{} times for 1{,}000 pockets (Figure~\ref{fig:cost}c).

The choice of the checkpoint matters more for the cross-attention network than for the Jev-like networks. The baseline networks of the preceding analysis choose their checkpoint on new targets, and with these checkpoints the dual encoder, the pooled network and the cross-attention network recover \kpDualTen{}, \kpPooledTen{} and \kpXattnTen{} on the 39 training pockets, against \kpConstTen{} for the constant pocket (Table~\ref{tab:known}; \kpDiffXattnDualTen{}, with the interval \kpDiffXattnDualTenLo{} to \kpDiffXattnDualTenHi{}, for the cross-attention network against the dual encoder). The three networks are then within \kpSpreadTen{} of each other, and the cross-attention network stops after \kpXaBestLo{} to \kpXaBestHi{} updates, against \kfXaBestLo{} to \kfXaBestHi{} updates when the checkpoint is chosen on known pockets.

\begin{table}[t]
\caption{Top-1\% recall of the networks on the 49 known pockets for ligands that were not used in training, within a budget of the best 1, 5, 10 and 20\% of the ranking (macro mean over the pockets; 95\% interval from bootstrapping the pockets). Each network is the mean over the cross-fitted folds that train on a pocket, with the checkpoint chosen on known pockets. The constant-pocket network is the dual encoder trained with one dummy pocket for every target. The pooled network replaces the product of the dual encoder by a learned readout. Time: seconds from SMILES to scores for the 260{,}155 ligands of DOCKSTRING against 1{,}000 pockets on one RTX 4090 (cross-attention in float16). The recall of every network with the checkpoint chosen on new targets and with longer training is in the Supporting Information.}
\label{tab:known}
\centering
\footnotesize
\setlength{\tabcolsep}{3pt}
\begin{tabular}{lcccrr}
\toprule
Network & 1\% & 5\% & 10\% (95\% interval) & 20\% & Time (s) \\
\midrule
Dual encoder & 0.29 & 0.72 & 0.88 (0.85 to 0.90) & 0.96 & 43.5 \\
Cross-attention & 0.35 & 0.79 & 0.92 (0.90 to 0.93) & 0.98 & 197 \\
Constant pocket & 0.25 & 0.64 & 0.81 (0.76 to 0.85) & 0.93 & 43.5 \\
Pooled concatenation (Jev-like, learned readout) & 0.31 & 0.75 & 0.89 (0.87 to 0.91) & 0.97 & 44.1 \\
\bottomrule
\end{tabular}
\end{table}

